\documentclass[11pt]{article}
\usepackage[margin=0.8in]{geometry}
\usepackage{amsmath,microtype}
\newcommand{\guideheading}[1]{\par\medskip\noindent{\large\bfseries #1}\par\smallskip}
\begin{document}
\begin{center}
{\Large Guide: Where the Current Merge Table Cannot Help}\par\medskip
Collatz proof project\quad---\quad September 5, 2026
\end{center}
\guideheading{The question}
The previous result let us combine several verified merge rules after
moving forward along an orbit. Could a fixed number of forward steps,
followed by as many rules as needed, cover every remaining candidate?
For the current table of 360 rules, the answer is no. There are systematic
arithmetic gaps, even when rule-chain length is unlimited.

\guideheading{The construction}
Choose any proposed forward horizon $K$. Take a seed $N$ satisfying
\[
N+1\text{ is divisible by }2^{K+18},\qquad
N\text{ is divisible by }3^{11}.
\]
These two requirements are compatible and give arbitrarily large seeds.
The first requirement makes the first $K$ steps odd. The second controls
which reverse rules can apply to the resulting values.

The exact residue check shows that at each of these forward values,
the only available rule undoes one odd step. Applying it repeatedly
returns to $N$, where no listed rule applies. Stopping earlier leaves
a value larger than $N$. No chain can get below the original seed.

\guideheading{What Lean verifies}
Lean checks twelve complete arithmetic classes. On eleven of them,
only the inverse odd rule can apply; on the initial class no rule can
apply. These claims hold for arbitrarily large members of the classes,
not just their smallest representatives. The finite checks use ordinary
kernel reduction and are included in the theorem audit.

The technical paper then proves, in ordinary mathematics, that the
construction above stays within those classes for any horizon $K$.
That arbitrary-horizon deduction has not yet been formalized in Lean.
The Python program independently reconstructs the classes and replays
sample horizons, including 128, with several progression quotients.

\guideheading{What this rules out}
Simply increasing a fixed forward-search limit and allowing longer
chains from the same table cannot establish universal coverage of
finite-prefix candidates. For every chosen limit, the construction
supplies candidates that defeat it. The seed depends on the limit;
this does not construct one orbit defeating every limit.

\guideheading{What remains possible}
The obstruction is specific to the current table and a uniform forward
horizon. It leaves open a proof that each candidate eventually reaches
a useful rule after a seed-dependent number of steps. New rules may also
break the obstruction. The constructed seeds have long noncontracting
prefixes, but they are not proved to be fully noncontracting or divergent.
This is a limitation of one proof approach, not a counterexample to Collatz.

\guideheading{Where to inspect the evidence}
\begingroup\raggedright
The formal module is \texttt{lean/Collatz/MergeRuleObstruction.lean}.
From \texttt{lean/}, run \texttt{lake build Collatz.MergeRuleObstruction}.
From the root, run \texttt{python3 analysis/merge\_rule\_obstruction.py
--horizon 64}. The adjacent test file checks residue agreement with Lean,
actual orbit steps, and the available rule outputs.
\par\endgroup
\end{document}
